Als Kreisradius nehmen wir die Drahtlänge, $r = \lhO$. Die Masse der Kugel spielt für die Kinematik keine Rolle. In einer Umlaufzeit $T = 1/f$ legt die Kugel den Kreisumfang $2\pi r$ zurück:
\begin{align}
 v &= \frac{2\pi r}{T} = \vhF \\
   &= 2\pi\cdot\lh\cdot\fh \\
   &= \vh \approx \result{\vhP{0}{2}}
\end{align}
Die Radialbeschleunigung ist
\begin{align}
 a_\text{r} &= \frac{v^2}{r} = \ahF \\
   &= 4\pi^2\cdot\lh\cdot\left(\fh\right)^2 \\
   &= \ah \approx \result{\ahP{0}{2}}
\end{align}

In Wirklichkeit ist der Kreisradius grösser als die Drahtlänge: Der Griff und die ausgestreckten Arme verlängern den Draht, und die Drehachse liegt im Körper des Athleten. Bei gleicher Frequenz wachsen $v \propto r$ und $a_\text{r} \propto r$, die Werte sind also \result{grösser}.
